\subsection{Coefficients of a Set of Classes of Modules}

Let C be a hereditary set of classes of A-modules (VIII, p.~183, Definition 1). We suppose that every A-module of type C is finite-dimensional over K and denote by \(\Theta_{\mathcal{C}}(\mathrm{A})\) the set of coefficients of the A-modules of type C. By Proposition 2 of VIII, p.~378, the set \(\Theta_{\mathcal{C}}(\mathrm{A})\) is also the union of the images of the A-linear mappings \(u: E \to A^{*}\) , where E runs through C; it is also the union of the subspaces \(\Theta_{\mathrm{E}}(\mathrm{A})\) of \(A^{*}\) , where E runs through C (loc. cit.). The family of (A, A)-sub-bimodules \((\Theta_{\mathrm{E}}(\mathrm{A}))_{\mathrm{E} \in \mathcal{C}}\) of \(A^{*}\) is directed, so its union \(\Theta_{\mathcal{C}}(\mathrm{A})\) is an (A, A)-sub-bimodule of \(A^{*}\) .

PROPOSITION 7. --- A left A-submodule of \(A^{*}\) of finite dimension over K is contained in \(\Theta_{\mathcal{C}}(\mathrm{A})\) if and only if it is of type C.

Let V be a left A-submodule of \(A^{*}\) of finite dimension over K. We saw in VIII, p.~379, Remark 3 that we have \(V \subset \Theta_{\mathrm{V}}(\mathrm{A})\) . If V is of type C, then we have \(\Theta_{\mathrm{V}}(\mathrm{A}) \subset \Theta_{\mathcal{C}}(\mathrm{A})\) , so V is contained in \(\Theta_{\mathcal{C}}(\mathrm{A})\) . Conversely, suppose that V is contained in \(\Theta_{\mathcal{C}}(\mathrm{A})\) . Since \(\Theta_{\mathcal{C}}(\mathrm{A})\) is the union of the directed family \((\Theta_{\mathrm{E}}(\mathrm{A}))_{\mathrm{E}\in\mathcal{C}}\) and V is finite-dimensional over K, there exists a module E of type C such that V is contained in \(\Theta_{\mathrm{E}}(\mathrm{A})\) . Now, there exists a natural number n such that \(\Theta_{\mathrm{E}}(\mathrm{A})\) is isomorphic to a quotient of \(\mathrm{E}^{n}\) (VIII, p.~378, Proposition 2). Since C is hereditary, V is of type C.

COROLLARY. --- Let E be an A-module of finite dimension over K. Then E is of type C if and only if \(\Theta_{\mathrm{E}}(\mathrm{A})\) is contained in \(\Theta_{\mathcal{C}}(\mathrm{A})\) .

The condition is obviously necessary.

Conversely, suppose that \(\Theta_{\mathcal{C}}(\mathrm{A})\) contains \(\Theta_{\mathrm{E}}(\mathrm{A})\) . The A-module \(\Theta_{\mathrm{E}}(\mathrm{A})\) is finite-dimensional over K (VIII, p.~378, Proposition 2), so it is of type C by Proposition 7. Now, there exists an integer \(n \geqslant 0\) such that E is isomorphic to an A-submodule of \(\Theta_{\mathrm{E}}(\mathrm{A})^{n}\) (VIII, p.~378, Proposition 2). Since C is hereditary, E is of type C.

